\chapter{Conclusão}
\label{cap:conclusao}

Esta tese apresentou uma arquitetura de navegação em duas camadas para inspeção de
subestações elétricas: uma patrulha deliberativa de sensoriamento ativo conduz os
sensores até onde eventos possam ocorrer, e uma camada reativa produz um impulso de
aproximação acionado por limiar em direção a anomalias e um impulso de recuo diante
de riscos, a partir de um campo de atração/repulsão $\Phi$ derivado da saliência
sensorial e arbitrado por uma precedência categórica de segurança. A arquitetura
roda sobre uma pilha madura e reprodutível (ROS~2 $+$ Nav2/MPPI $+$ Gazebo Fortress).

\section{Síntese das Atribuições}
\label{sec:conc-atribuicoes}

A avaliação do Capítulo~\ref{cap:resultados} \emph{estreita} aquilo que pode ser
creditado à camada reativa. As três atribuições são enunciadas aqui sem rodeios,
porque a precisão quanto ao que cada mecanismo entrega é parte da contribuição.

\subsection{O Ganho de Segurança no Risco Estático Vem da Zona de Exclusão}

Sobre um risco \emph{estático} colocado na rota de patrulha, uma patrulha fixa
portando a mesma zona de exclusão disparada por detecção e nenhum supervisor reativo
(KOF) e a ablação somente-zona que retém a inspeção reativa (KOa) alcançaram
afastamentos mínimos medianos semelhantes ($1{,}73$ e $1{,}71$~m), sem diferença
estatisticamente detectável ($p_{\mathrm{Holm}}=0{,}4194$); o método completo (RHM)
situa-se nos mesmos $1{,}71$~m.

%% Delimitado em 13/11/2026. A atribuicao vale para o cenario e as configuracoes
%% estudadas; ausencia de beneficio adicional detectado nao demonstra equivalencia
%% geral entre os componentes.
Nesse cenário e nessas configurações, portanto, o benefício de segurança é atribuível
à zona de exclusão deliberativa, e não à ponderação reativa. A afirmação é dessa
extensão: ela não estabelece equivalência geral entre os dois componentes, nem se
transfere a riscos que não estejam estáticos sobre a rota: a comparação de risco
emergente da Seção~\ref{sec:res-emergente} mostra justamente o caso em que a resposta
reativa acrescenta margem. O que a camada reativa acrescenta sobre o KOF são as
inspeções em si ($3$ contra $0$ de $3$), ao custo de cerca de $31$~s de tempo de
missão adicional para o KOa.

\subsection{O Critério de Seleção Não Teve Efeito Detectável; o Modo de Disparo Teve}

Uma variante que difere apenas por selecionar o alvo mais próximo em vez do mais
saliente produziu os mesmos desfechos discretos constantes que a camada proposta ao
longo das dez execuções nos cenários \texttt{default} e \texttt{conflict}, e
distribuições sobrepostas em \texttt{adverse}. Não se alega equivalência: apenas
que nenhuma diferença foi detectada neste tamanho amostral.

O que \emph{mudou} o desfecho foi o \emph{modo de disparo}: uma heurística genérica
de detectar--selecionar--desviar--retornar, sobre o mesmo controlador Nav2/MPPI e
acionada por detecção bruta, acaba por inspecionar a anomalia colocalizada com um
risco e reduz o afastamento mínimo à metade ($4{,}84$ contra $2{,}23$~m, $\delta$ de
Cliff $-1{,}00$ na convenção comparador-menos-método adotada em toda a tese: o braço
comparado é o menor), ao mesmo tempo em que atrasa a inspeção prioritária em cerca de
$105$~s.

Essa troca \emph{não} remove a precedência categórica --- a heurística ainda
preempta, três vezes por execução ---, mas altera conjuntamente o sinal de decisão e
sua persistência temporal, de modo que os dois não são separáveis com as $32$
configurações da matriz principal. Os \textit{logs} de episódio são consistentes com
uma explicação ligada à persistência, e uma verificação posterior a sustenta em parte:
encurtar o período de decaimento \emph{mantendo fixa} a fonte do disparo encurta o
tempo em investigação com separação completa no braço que dispara sobre o campo, e não
separa nenhum desfecho no braço que dispara por detecção bruta
(§\ref{sec:res-relogio-consistente}).

%% CORRIGIDO em 16/11/2026. "Nao move desfecho algum" afirmava efeito nulo; o que houve
%% foi ausencia de separacao completa, com deltas pequenos e nao nulos, e sem teste de
%% equivalencia que autorizasse a leitura forte.
Não separar não é o mesmo que não mover: os deltas naquele braço são pequenos, não
nulos, e o protocolo não inclui teste de equivalência. A persistência, portanto, não é
inerte onde é consultada; sobre o braço que não a consulta, o que se tem é ausência de
separação detectada. Que fração do efeito no cenário de conflito lhe cabe permanece sem medida,
porque essa verificação encurta a persistência em vez de removê-la e correu no cenário
\texttt{default}.

\subsection{A Precedência Tem Custo Medido}

%% Delimitado em 15/11/2026. A frase anterior juntava num so enunciado a
%% janela de ~70 s e a inspecao perdida. Sob o relogio especificado a inspecao perdida
%% persiste (10/10 tambem) e a janela nao: o recuo que a abre deixa de disparar. Os
%% dois fatos foram separados porque tem estatutos diferentes.
Essa precedência carrega um custo que foi medido: um risco marginalmente elegível faz o
robô abrir mão de uma de quatro inspeções em $10/10$ execuções, e isso se repete nas
duas implementações temporais testadas: $20/20$ ao todo. A supressão da atração por
cerca de $70$~s que acompanha esse custo, e o recuo único que a abre, pertencem ao
relógio efetivamente executado: sob o período especificado o recuo não dispara
(§\ref{sec:res-relogio-consistente}). O que se repete nas duas implementações é o
\emph{desfecho}, três de quatro inspeções, e não a sua causa: sob o relógio especificado
duas das investigações terminaram por aborto de navegação
(§\ref{sec:res-relogio-consistente}), de modo que a mesma perda não é atribuída ao mesmo
mecanismo nos dois casos. Resultado repetido entre condições não identifica
mecanismo.

%% Delimitado em 13/11/2026. A frase anterior afirmava sem condicao que o
%% distintivo e o canal de risco "e nao a ponderacao por saliencia". A segunda
%% metade e ausencia de efeito DETECTADO nos cenarios e no tamanho amostral
%% executados, e nao demonstracao de que a ponderacao seja inutil.
Em conjunto, nos cenários executados e com oito a dez execuções por braço, o que
distingue esta camada de uma heurística genérica de desviar e retornar é o canal de
risco com limiar e decaimento. A ponderação por saliência, no mesmo conjunto de
condições, não produziu diferença detectável, o que delimita a sua contribuição
\emph{nestas distribuições de eventos} e não estabelece que ela seja dispensável em
outras, com anomalias de intensidades mais separadas ou em maior número. E a camada
troca completude por afastamento de um modo que está agora quantificado nas duas
direções.

\section{Limites da Evidência}
\label{sec:conc-limites}

%% Reescrito em 13/11/2026: os numeros da reavaliacao do limiar de
%% confirmacao apareciam aqui pela primeira vez, o que fazia da conclusao um capitulo
%% com resultado inedito. Eles passam a ser apresentados na
%% Secao~\ref{sec:res-confirmacao}, e aqui apenas se sintetiza.
A completude de inspeção é relatada como a taxa de execuções que inspecionam
\emph{todas} as anomalias: $10/10$ com percepção limpa, $4/9$ sob ruído alto de
percepção e $9/10$ sob aglomeração densa. Conforme a
Seção~\ref{sec:res-confirmacao}, essa taxa mistura dois efeitos que a análise separa e
que ambos existem: o critério de confirmação é sensível ao limiar adotado, e a
aproximação física também se degrada sob ruído. Elevar o limiar não corrige a segunda
coisa --- apenas deixa de contar a primeira ---, razão pela qual o critério adotado
permanece o justificado por construção.

Nenhuma entrada na região de contato de $0{,}5$~m foi observada em qualquer
configuração, mas com $8$ a $10$ execuções válidas por grupo o limite superior
unilateral exato de $95\%$ sobre a probabilidade por execução de tal entrada
permanece entre $25{,}9\%$ e $31{,}2\%$; e, sobre o risco na rota, o método transita
\emph{dentro} da faixa consultiva de afastamento de $2$~m que ele próprio relata, por
cerca de $9$~s por execução. Esses são os limites da evidência de segurança, e não
uma garantia de segurança.

\section{Escopo}
\label{sec:conc-escopo}

A avaliação é em simulação. A percepção é abstraída em eventos-oráculo, e não em
fluxos brutos de sensores, e a transformada de localização é mantida fixa. Os
resultados caracterizam, portanto, a \emph{lógica de decisão} da arquitetura e não
devem ser lidos como validação de um sistema de inspeção autônoma pronto para
implantação.

A região de contato de $0{,}5$~m e a faixa de afastamento de $2$~m são distâncias
centro a centro, definidas operacionalmente como aproximações de medida --- com raio
do robô de $0{,}25$~m ---, e não sensoriamento físico de contato nem distâncias de
isolamento elétrico normativas.

%% Movido para ca em 12/11/2026. O paragrafo estava dentro de \section{Limitações},
%% o que enquadrava uma decisao de projeto como limitacao. A justificativa completa,
%% com os achados que a motivaram e a tabela objetivo-por-objetivo, esta na
%% Secao~\ref{sec:intro-escopo}.
Esses limites decorrem das decisões de delimitação tomadas durante o
desenvolvimento e justificadas na Seção~\ref{sec:intro-escopo}: a abstração da
percepção em eventos e a fixação da transformada de localização foram adotadas como
controles de confundimento, para que o efeito de cada mecanismo da arquitetura
pudesse ser atribuído. A proposta de qualificação previa, em seu lugar, um
\textit{framework} de cinco módulos integrando fusão sensorial heterogênea,
aprendizado por reforço e visão computacional; os resultados reunidos ao longo desta
tese --- a atribuição do afastamento à zona de exclusão e não à ponderação reativa, e
a constatação de que a tabela de valor não seleciona as ações executadas ---
redirecionaram o trabalho para a arquitetura focada de duas camadas aqui descrita. A
fundamentação dos Capítulos~\ref{cap:fundamentacao} e~\ref{cap:revisao} preserva o
contexto mais amplo que motivou a investigação; os
Capítulos~\ref{cap:proposta} a~\ref{cap:resultados} descrevem o que foi entregue, e
a Seção~\ref{sec:conc-futuro} nomeia o que não foi.

\section{Limitações}
\label{sec:conc-limitacoes}

Erro de localização, persistência de eventos, falsos positivos e negativos e
geometria de mapa permanecem não modelados. O cenário adverso é uma única
configuração densa e a varredura de ruído cobre dois níveis, de sorte que nenhum dos
dois sustenta alegação geral de escalabilidade.

%% Acrescentado em 15/11/2026: a limitacao de implementacao mais consequente
%% faltava nesta lista, que so trazia limitacoes de modelagem.
Uma limitação de \emph{implementação} acompanha essas de modelagem, e é declarada com
o mesmo peso. A guarda que autoriza o decaimento dos acumuladores compara instantes de
relógio de parede enquanto os temporizadores seguem o relógio simulado, de modo que o
período realizado nas $320$ execuções analisadas foi de cerca de $1{,}6$~s onde a
especificação pede $1$~s (Seção~\ref{sec:prop-relogios}). Toda conclusão que dependa de
\emph{quanto tempo} a memória persiste é, portanto, afirmada da implementação avaliada.
Três configurações foram reexecutadas sob o relógio especificado para medir o alcance
dessa dependência (Seção~\ref{sec:res-relogio-consistente}); as outras vinte e nove
não.

%% Acrescentado em 13/11/2026: o que as repeticoes licenciam.
As dez repetições por configuração estabelecem \emph{repetibilidade} daquele arranjo
sob a não-determinação do simulador, e não mais do que isso: fora das duas
configurações de ruído não há aleatorização planejada entre execuções
(Seção~\ref{sec:res-ameacas}). Repetir o mesmo \textit{layout} e a mesma distribuição
de eventos dez vezes não equivale a avaliar dez instalações, e os cenários não são
amostra de população alguma: são arranjos escolhidos para exercitar mecanismos
específicos. Daí também que diferenças pequenas possam exibir separação estatística
completa sem relevância operacional equivalente, razão pela qual as magnitudes são
relatadas ao lado dos valores de $p$.

A tabela de valor é tabular e, conforme a Seção~\ref{sec:prop-atualizacao}, não
seleciona as ações executadas; não se faz, portanto, qualquer alegação de
convergência ou de invariância de política para a atualização aumentada. O que a
Tabela~\ref{tab:res-principal} atesta, pelo resíduo completo
$|\alpha\delta+\eta\Phi|$, é que os incrementos permanecem limitados nas simulações
avaliadas: estabilidade dos incrementos, e não convergência.

\section{Trabalhos Futuros}
\label{sec:conc-futuro}

%% Corrigido em 14/11/2026: a abertura excluia a ponderacao como lugar de
%% interesse de forma categorica. Nao detectar efeito em tres cenarios nao elimina a
%% utilidade da saliencia em outras distribuicoes de eventos -- afirmar isso e
%% justamente o tipo de generalizacao que o resto do capitulo recusa.
A ausência de efeito detectável da regra de seleção e o custo quantificado da
persistência apontam na mesma direção: \emph{nas distribuições de eventos estudadas},
a arbitragem e sua extensão temporal são onde o trabalho adicional rende mais do que a
ponderação. A delimitação importa: os cenários executados colocam poucos eventos
simultâneos e bem separados, condições em que saliência e proximidade tendem a
coincidir na escolha. Em distribuições mais densas, ou com eventos de intensidades
muito desiguais a distâncias comparáveis, a ponderação pode voltar a pesar, o que
este trabalho não testou e, portanto, não exclui.

\begin{itemize}
%% Reescrito em 13/11/2026: a comparacao contra Q-learning tabular, que este item
    %% dava como pendente, foi executada (Secao~\ref{sec:res-qdriven}). O que resta em
    %% aberto e a politica de DEFERENCIA, treinada no proprio ambiente de experimentacao.
    \item \textbf{Política aprendida de deferência.} Uma política que
    \emph{aprenda} quando deferir ao canal de risco, em vez de sempre deferir,
    endereçaria o excesso de conservadorismo que foi medido. A formulação de valor da
    Seção~\ref{sec:prop-atualizacao} é o substrato natural para isso. A comparação
    contra uma política tabular treinada com recompensa aumentada já foi conduzida
    (Seção~\ref{sec:res-qdriven}), com piora do braço aprendido no tempo em investigação
    e na missão adversa e resultados mistos nas demais métricas, sem ordenação global
    entre os dois; resta, porém, a
    pergunta que ela não responde: uma política treinada no \emph{próprio}
    ambiente de avaliação, com ações mapeadas para o controlador em vez de para a
    velocidade do robô, e com a deferência ao canal de risco como objeto do
    aprendizado em vez de regra fixa.

%% Acrescentados em 13/11/2026: os dois desenhos que permitiriam alegacoes isoladas
    %% onde hoje so ha o conjunto de intervencoes. Ficam nomeados aqui para que o
    %% limite das atribuicoes dos Capitulos~\ref{cap:resultados} e~\ref{cap:conclusao}
    %% tenha o experimento correspondente identificado.
%% Reescrito em 15/11/2026: o item dava o desenho como integralmente
    %% ausente. A bateria de relogios executou a sua primeira metade -- variar a
    %% persistencia com o disparo fixo. O que falta e a outra metade e a replicacao no
    %% cenario de conflito.
    \item \textbf{Separar disparo de persistência.} Um desenho que varie
    independentemente a fonte do disparo e a extensão temporal da memória, mantendo
    constantes a seleção, a prioridade e o controlador, isolaria o fator que a matriz
    principal só atribui ao par. Metade desse desenho já correu: variar a persistência
    com o disparo fixo separou um desfecho no braço que consulta o campo e não separou
    nenhum no que dispara por detecção bruta, sem que daí se conclua efeito nulo neste
    (Seção~\ref{sec:res-relogio-consistente}). Faltam a
    outra metade --- variar a fonte do disparo com a persistência fixa --- e a
    replicação no cenário de conflito, onde o efeito a decompor foi medido.

%% Acrescentado em 15/11/2026.
    \item \textbf{Reexecutar a matriz sob o relógio especificado.} Três das $32$
    configurações foram verificadas com a guarda do decaimento alimentada pelo relógio
    do temporizador; nelas os contrastes centrais sobreviveram e a janela de supressão
    não. Estender essa verificação às demais trocaria a transposição declarada na
    Seção~\ref{sec:res-ameacas} por medida direta. Fica como trabalho futuro, e não
    como correção pendente, por delimitação de escopo: uma bateria nova pode coexistir
    com a histórica, como o lote de verificação já faz, e o que o Capítulo~\ref{cap:resultados}
    afirma está declarado como sendo da implementação avaliada.

    \item \textbf{Separar interrupção de deslocamento.} Uma variante de parada e
    espera --- mesma preempção da investigação, mesmo tratamento do sinal de risco,
    sem a meta de recuo --- separaria o efeito de interromper a aproximação do efeito
    de deslocar-se, que a comparação de risco emergente hoje mede em conjunto.

    \item \textbf{Percepção real.} Substituir os eventos-oráculo por percepção
    térmica, acústica e visual real alimentando a camada de eventos, com tratamento
    explícito de falsos positivos e negativos.

    \item \textbf{Localização realista.} Liberar a transformada
    \texttt{map}$\rightarrow$\texttt{odom} e medir o efeito da deriva de localização
    sobre a confirmação de inspeção e sobre as métricas de afastamento.

    \item \textbf{Validação em \textit{hardware}.} Validação em \textit{hardware
    in the loop} sobre a plataforma quadrúpede GO2~PRO e ensaio de campo em
    subestação real.

    \item \textbf{Varreduras mais amplas.} Densidade de eventos, geometria de mapa e
    níveis de ruído adicionais, para sustentar alegações de escalabilidade que uma
    única configuração adversa não sustenta.
\end{itemize}

%% Secao criada em 13/11/2026. O quadro abaixo permite conferir cada
%% afirmacao desta tese contra a evidencia que a sustenta, as condicoes em que foi
%% obtida e o limite que ela nao ultrapassa, sem percorrer o capitulo de resultados.
\section{Quadro de Afirmações}
\label{sec:conc-quadro}

A \autoref{tab:conc-quadro} reúne cada afirmação desta tese com a evidência que a
sustenta, as condições em que foi obtida e o limite que ela não ultrapassa. Nenhuma
linha contém resultado que não esteja no Capítulo~\ref{cap:resultados}.

\begin{table}[htb]
\centering
\caption{Afirmação, evidência, condições e limite. Os ponteiros remetem às seções de
resultados correspondentes.}
\label{tab:conc-quadro}
\scriptsize
\setlength{\tabcolsep}{3pt}
\begin{tabular}{p{0.215\textwidth}p{0.215\textwidth}p{0.195\textwidth}p{0.225\textwidth}}
\hline
\textbf{Afirmação} & \textbf{Evidência} & \textbf{Condições} & \textbf{Limite}\\
\hline
O afastamento diante de risco estático na rota vem da zona de exclusão deliberativa &
KOF e KOa com afastamentos medianos de $1{,}73$ e $1{,}71$~m, sem diferença detectável
($p_{\mathrm{Holm}}=0{,}4194$); §\ref{sec:res-ablacao} &
Cenário de risco estático sobre a rota; $8$--$10$ execuções por braço &
Não demonstra equivalência geral de componentes; não se transfere a risco emergente\\
\hline
Sem a zona de exclusão, o recuo reativo sozinho não resolve o risco na rota &
Nenhuma das nove trajetórias válidas de KOb concluiu a missão no orçamento; uma décima
tentativa permaneceu imóvel e foi excluída; interrupção em $627\pm6$~s;
§\ref{sec:res-ablacao} &
Mesmo cenário; orçamento de parede de $600$~s &
$627\pm6$~s são instantes de interrupção, não tempos de conclusão; a duração de conclusão
é censurada à direita e não tem limiar comum entre trajetórias\\
\hline
A camada reativa acrescenta as inspeções que a patrulha com zona de exclusão não faz &
$3$ de $3$ contra $0$ de $3$; custo de $\approx 31$~s; §\ref{sec:res-ablacao} &
Mesmo cenário &
Inspeção é proxy de chegada creditada (§\ref{sec:setup-metricas})\\
\hline
O critério de seleção não teve efeito detectável &
Desfechos discretos idênticos em dez execuções; §\ref{sec:res-heu} &
Três cenários; sem ruído de percepção &
Ausência de efeito detectado; nenhum teste de equivalência aplicado a este contraste\\
\hline
O modo de disparo altera o desfecho de segurança &
$4{,}84$ contra $2{,}23$~m de afastamento mínimo, $\delta=-1{,}00$; atraso de
$\approx 105$~s; §\ref{sec:res-heu} &
Cenário de conflito, risco colocalizado com anomalia mais saliente &
Disparo e persistência variam em conjunto; não são fatores separáveis
(§\ref{sec:res-ameacas})\\
\hline
Em CM, as duas implementações temporais terminaram com $3$ de $4$ inspeções &
Uma de quatro inspeções abandonada em $10/10$ sob cada relógio ($20/20$);
§\ref{sec:res-custo-precedencia} &
Cenário de risco marginalmente elegível, nas duas implementações temporais testadas
(§\ref{sec:res-relogio-consistente}) &
Mecanismos observados e \emph{não} isolados: sob o relógio de parede, supressão da
atração após o recuo; sob o especificado, aborto de navegação em duas investigações. O
desfecho se repete; a causa não foi identificada\\
\hline
A supressão de $\approx 70$~s e o recuo que a abre são da implementação &
Janela de $69{,}5$~s e um recuo sob o relógio executado; $0{,}00$ recuos
($\delta=-1{,}00$) sob o especificado; §\ref{sec:res-relogio-consistente} &
Configuração marginal, $10$ execuções por relógio &
Afirmação sobre o código avaliado, \emph{não} sobre a precedência categórica em si\\
\hline
A preempção e o recuo compram margem contra risco emergente &
DHa supera DHb em $0{,}33$~m e $0{,}51$~s, com separação completa ($|\delta|=1$) no
afastamento; §\ref{sec:res-emergente} &
Risco que ativa a $3$~m durante investigação; $8$ execuções válidas &
Efeito do \emph{conjunto} de ações, não do deslocamento isolado; métricas de missão
inteira\\
\hline
Controlar o movimento por tabela de valor não melhorou a subtarefa reativa &
Mesma completude; tempo em investigação maior com $\delta=+1{,}00$ nos três cenários;
demais métricas mistas; §\ref{sec:res-qdriven} &
Três cenários; política treinada em simulador cinemático &
Comparação entre as duas implementações testadas; camada de eventos e controlador não
compartilhados (§\ref{sec:setup-qdriven}); não testa arbitragem aprendida entre risco
e inspeção nem isola o efeito do aprendizado\\
\hline
A completude se degrada sob ruído alto de percepção &
$4/9$ execuções completas; seis rejeições na confirmação;
§\ref{sec:res-ruido} &
Duas realizações de ruído; $9$ execuções válidas &
Mistura sensibilidade do critério e piora da aproximação, ambas quantificadas
(§\ref{sec:res-confirmacao})\\
\hline
Nenhuma entrada na região de contato foi observada &
Zero entradas em todas as configurações; §\ref{sec:res-seguranca} &
$8$--$10$ execuções por grupo &
Limite superior de $95\%$ entre $25{,}9\%$ e $31{,}2\%$; não é colisão física medida\\
\hline
Os incrementos da atualização permanecem limitados &
Resíduo completo $|\alpha\delta+\eta\Phi|$ tabulado;
Tabela~\ref{tab:res-principal} &
Todas as configurações com atualização ativa &
Estabilidade dos incrementos, não convergência; a tabela não governa o comportamento\\
\hline
\end{tabular}
\end{table}

\section{Trajetória de Transferência Tecnológica}
\label{sec:conc-transferencia}

As decisões que esta pesquisa permite informar são as mensuráveis, e só elas. Com os
números acima, um projeto de inspeção robótica pode estimar o custo de missão por
investigação adicional, o efeito da memória de risco sobre reativações, a perda de
completude imposta por uma precedência de segurança rígida e as situações em que a
expressão deliberativa do risco já responde pelo afastamento, tornando dispensável
investir em ponderação reativa para aquele caso. São parâmetros de projeto de
supervisão, não garantias de planejamento de ronda: cada valor está condicionado ao
cenário em que foi obtido, conforme a \autoref{tab:conc-quadro}.

O que \emph{não} se oferece: estimativa financeira, nível de maturidade tecnológica,
tempo de implantação ou ganho operacional em instalação real. Nada nesta tese sustenta
tais números, e convertê-los a partir de resultados de simulação seria extrapolação.

A trajetória de transferência, portanto, é formada por etapas \emph{não executadas},
que se nomeiam como tais e em ordem de dependência: (i) substituir os eventos-oráculo
por percepção térmica, acústica e visual real, com tratamento de falsos positivos e
negativos, único passo que torna o restante mensurável em campo; (ii) liberar a
transformada de localização e medir o efeito da deriva sobre confirmação de inspeção e
afastamento; (iii) validação em \textit{hardware in the loop} sobre a plataforma
quadrúpede, onde dinâmica e atraso de atuação entram em cena; e (iv) ensaio em
subestação real, sob os requisitos normativos e de segurança da instalação, que esta
tese não trata. Cada etapa é condição da seguinte, e nenhuma delas foi realizada aqui.

\section{Considerações Finais}
\label{sec:conc-finais}

A contribuição desta tese não é a demonstração de que um robô pode inspecionar uma
subestação --- isso já se sabia --- nem a de que um campo de atração e repulsão pode
conduzir um desvio, o que uma heurística simples também faz. A contribuição é a
\emph{decomposição}: mostrar, com ablações sobre um controlador compartilhado e com
análise estatística que declara suas próprias fronteiras, qual mecanismo entrega qual
efeito. A zona de exclusão deliberativa compra o afastamento; a camada reativa compra
as inspeções; o canal de risco com limiar e decaimento compra a persistência da
precedência de segurança, e cobra por ela, em inspeções perdidas, um preço que
este trabalho mediu em vez de supor inexistente.

Em um domínio no qual alegações de autonomia costumam exceder a evidência que as
sustenta, delimitar o que de fato foi demonstrado parece contribuição tão legítima
quanto o próprio mecanismo.
